% TOP_PROBLEM: {"id": 30007380, "title": "A GFF limit for triangular odd heights with two boundary arcs", "category_id": 19, "category": {"id": 19, "name": "probability", "display_name": "Probability", "description": "Problems involving probability theory, stochastic processes, and random structures.", "slug": "probability", "order_index": 19, "created_at": "2026-07-31T15:26:25.671Z"}, "status": "open"}
% FIRST_POSTED: 2026-10-06
% ATTEMPT_STATUS: solved
% OPENAI_CLAIM: {"schema_version":1,"source_repo":"https://github.com/openai/math","source_commit":"adc7f1241b42e322a6451854ab7e4b4c146bf78a","release_date":"2026-10-06","match":"full","resolution":"proved","independently_reviewed":false,"second_pass":"agreed","families":[{"family":"232","title":"Gaussian fields and interfaces for triangular-lattice Lipschitz heights","subject":"Probability and statistical mechanics","manuscripts":[{"title":"The Gaussian free field limit of integer Lipschitz heights with two-arc boundary data","pdf_path":"preprints/The-Gaussian-free-field-limit-of-integer-Lipschitz-heights-with-two-arc-boundary-data-September-25-2026/paper.pdf","readme_path":"preprints/The-Gaussian-free-field-limit-of-integer-Lipschitz-heights-with-two-arc-boundary-data-September-25-2026/README.md","date":"2026-09-25","theorem_ref":"thm:main (sections/01-introduction.tex:47)","statement_sources":[{"path":"preprints/The-Gaussian-free-field-limit-of-integer-Lipschitz-heights-with-two-arc-boundary-data-September-25-2026/build/sections/01-introduction.tex","line":47,"label":"thm:main"}]}],"lean":null,"reasoning_trace":null}]}
\documentclass[11pt]{article}
\usepackage[margin=25mm]{geometry}
\usepackage{fontspec}
\setmainfont{Latin Modern Roman}
\usepackage{amsmath,amssymb,mathtools}
\usepackage{unicode-math}
\setmathfont{Latin Modern Math}
\usepackage{enumitem}
\usepackage{xurl,hyperref,bookmark}
\hypersetup{colorlinks=true,urlcolor=blue}
\setcounter{secnumdepth}{0}
\setlength{\parindent}{0pt}
\setlength{\parskip}{5pt}
\setlength{\emergencystretch}{3em}
\begin{document}
\section*{A GFF limit for triangular odd heights with two boundary arcs}

\subsection{OpenAI's claim}
OpenAI presents a proof in the selected manuscript. Two-arc odd integer Lipschitz heights on triangular lattices are claimed to converge, after centering and scalar normalization, to the Dirichlet Gaussian free field.

\subsection{Mathematical statement of the claimed result}
\paragraph{Result 1.} On triangular-lattice polygonal approximations $D_\delta$ of a bounded smooth simply connected domain $D$, take two distinct boundary marks at edge midpoints. Assume uniform convergence of marked boundary parametrizations. Sample odd integer vertex heights uniformly subject to
\[
 h_\delta=1\text{ on }A_{+,\delta},\quad h_\delta=-1\text{ on }A_{-,\delta},\quad
 |h_\delta(x)-h_\delta(y)|\in\{0,2\}\ (x\sim y).
\]
Interpolate affinely over triangles. If $\bar g$ is harmonic with boundary values $+1,-1$ on the respective limiting arcs and $\Phi_D$ is the Dirichlet GFF with covariance $(-\Delta_D)^{-1}$, a lattice-dependent $\sigma>0$ satisfies
\[
 h_\delta\Longrightarrow\bar g+\sigma\Phi_D\quad\text{in }\mathcal D'(D),\qquad
 \mathbb Eh_\delta(f)\to\int_D\bar gf\quad(f\in C_c^\infty(D)).
\]
Consequently $(h_\delta-\mathbb Eh_\delta)/\sigma\Longrightarrow\Phi_D$. All finite joint test-function laws and their mixed moments converge. The normalization is $\sigma=4\pi\sqrt{c_*}$ for the paper's absolutely convergent finite-volume coefficient $c_*$.

\subsection{Outline of the argument}
Exact height cuts are combined with buffered circuit mixing and rare-pin estimates. Reflection positivity fixes the plane covariance and an absolutely convergent normalization; moving boundary pins and conditional Gaussian moment identities identify the field in the prescribed domain.

\subsection{Scope relative to this problem}
Two-arc odd integer Lipschitz heights on triangular lattices are claimed to converge, after centering and scalar normalization, to the Dirichlet Gaussian free field. The question retains these same parameters and conventions. Within these assumptions, this is a claimed affirmative answer to the stated question.

\subsection{Formal verification}
Family 232 has no Lean scope page in this release. Its selected result is not presented as formalized. This site has not executed the Lean code. The scope descriptions do not certify the remaining manuscript arguments.

\subsection{Status caveat}
OpenAI attributes these manuscripts to work by an internal model. Its release cautions that arguments without formalization can contain mistakes. This catalogue has not independently checked the proofs; the labels report the mathematical claims attributed to OpenAI.

\subsection{Sources}
\begin{itemize}
\item OpenAI, \emph{The Gaussian free field limit of integer Lipschitz heights with two-arc boundary data}, 2026-09-25; citation key \texttt{OAI:The-Gaussian-free-field-limit-of-integer-Lipschitz-heights-with-two-arc-boundary-data-September-25-2026}. \href{https://github.com/openai/math/blob/adc7f1241b42e322a6451854ab7e4b4c146bf78a/preprints/The-Gaussian-free-field-limit-of-integer-Lipschitz-heights-with-two-arc-boundary-data-September-25-2026/paper.pdf}{PDF (GitHub)}; \href{https://github.com/openai/math/raw/adc7f1241b42e322a6451854ab7e4b4c146bf78a/preprints/The-Gaussian-free-field-limit-of-integer-Lipschitz-heights-with-two-arc-boundary-data-September-25-2026/paper.pdf}{PDF (direct)}; \href{https://github.com/openai/math/blob/adc7f1241b42e322a6451854ab7e4b4c146bf78a/preprints/The-Gaussian-free-field-limit-of-integer-Lipschitz-heights-with-two-arc-boundary-data-September-25-2026/README.md}{Manuscript page and citation}; \href{https://github.com/openai/math/blob/main/preprints/The-Gaussian-free-field-limit-of-integer-Lipschitz-heights-with-two-arc-boundary-data-September-25-2026/paper.pdf}{Latest version}; Statement source: \href{https://github.com/openai/math/blob/adc7f1241b42e322a6451854ab7e4b4c146bf78a/preprints/The-Gaussian-free-field-limit-of-integer-Lipschitz-heights-with-two-arc-boundary-data-September-25-2026/build/sections/01-introduction.tex\#L47}{thm:main}.
\item \href{https://github.com/openai/math/blob/adc7f1241b42e322a6451854ab7e4b4c146bf78a/CONTENTS.md}{OpenAI catalogue entry}.
\end{itemize}
\end{document}
